\documentclass[11pt]{article}
\usepackage[utf8]{inputenc}
\usepackage[spanish, es-tabla]{babel}
\usepackage{amsmath, amssymb}
\usepackage{booktabs}
\usepackage{enumitem}
\usepackage[margin=2.5cm]{geometry}

\setlength{\parindent}{0pt}
\setlength{\parskip}{0.65em}

\title{Econometría ICS-294\\Guía de ejercicios: Clase 6}
\author{Francisco Alfaro Medina}
\date{}

\begin{document}
\maketitle

\section*{Objetivo}

Practicar la estimación del modelo de regresión lineal múltiple, su interpretación ceteris paribus, la notación matricial, las propiedades algebraicas de MCO y las medidas de bondad de ajuste.

\section*{Ejercicios}

\begin{enumerate}[label=\textbf{Ejercicio \arabic*.}, leftmargin=*]

\item \textbf{Interpretación ceteris paribus.}

Considere el modelo:
\[
salario_i=\beta_0+\beta_1 educ_i+\beta_2 exper_i+\beta_3 mujer_i+u_i.
\]

\begin{enumerate}[label=\alph*)]
    \item Interprete \(\beta_1\), \(\beta_2\) y \(\beta_3\).
    \item Explique por qué la interpretación de \(\beta_1\) es ceteris paribus.
    \item Mencione dos variables que podrían estar omitidas en \(u_i\).
\end{enumerate}

\item \textbf{Predicción con múltiples regresores.}

Suponga que se estima:
\[
\widehat{salario}=100+40\,educ+10\,exper-30\,mujer.
\]

\begin{enumerate}[label=\alph*)]
    \item Prediga el salario de un hombre con 12 años de educación y 5 años de experiencia.
    \item Prediga el salario de una mujer con las mismas características.
    \item Interprete la diferencia entre ambas predicciones.
\end{enumerate}

\item \textbf{Forma matricial.}

Para el modelo \(y_i=\beta_0+\beta_1x_{1i}+\beta_2x_{2i}+u_i\), considere:
\[
X=\begin{bmatrix}
1 & 1 & 2\\
1 & 2 & 1\\
1 & 3 & 4
\end{bmatrix},
\qquad
y=\begin{bmatrix}
5\\
6\\
10
\end{bmatrix}.
\]

\begin{enumerate}[label=\alph*)]
    \item Identifique \(n\) y \(k\).
    \item Escriba el vector \(\beta\).
    \item Calcule \(X^\top y\).
    \item Explique qué representa \(X^\top X\).
\end{enumerate}

\item \textbf{Condiciones normales.}

Explique por qué las condiciones de primer orden de MCO pueden escribirse como:
\[
X^\top \hat{u}=0.
\]

Luego interprete esta condición para:
\begin{enumerate}[label=\alph*)]
    \item la constante;
    \item una variable explicativa \(x_1\);
    \item todas las variables explicativas en conjunto.
\end{enumerate}

\item \textbf{Propiedades algebraicas.}

En una regresión con constante, explique por qué se cumple:
\[
\sum_{i=1}^{n}\hat{u}_i=0,
\qquad
\bar{y}=\bar{\hat{y}}.
\]

Use estas propiedades para explicar por qué la regresión estimada pasa por el punto de medias.

\item \textbf{\(R^2\) y \(R^2\) ajustado.}

Para un modelo con \(n=80\), \(k=3\), \(SSR=240\) y \(SST=1000\):

\begin{enumerate}[label=\alph*)]
    \item Calcule \(R^2\).
    \item Calcule \(\overline{R^2}\).
    \item Interprete ambos resultados.
    \item Explique por qué \(\overline{R^2}\) puede ser menor que \(R^2\).
\end{enumerate}

\item \textbf{Comparación de modelos.}

Dos modelos estimados con la misma variable dependiente entregan:

\begin{center}
\begin{tabular}{lccc}
\toprule
Modelo & \(n\) & \(k\) & \(SSR\) \\
\midrule
A & 100 & 2 & 500 \\
B & 100 & 5 & 460 \\
\bottomrule
\end{tabular}
\end{center}

Suponga que \(SST=1200\).

\begin{enumerate}[label=\alph*)]
    \item Calcule \(R^2\) para cada modelo.
    \item Calcule \(\overline{R^2}\) para cada modelo.
    \item ¿Cuál modelo preferiría si se quiere penalizar la inclusión de regresores adicionales?
\end{enumerate}

\item \textbf{Multicolinealidad perfecta.}

Suponga que en una base de datos se incluyen \(educ\), \(educ\_meses=12\cdot educ\) y una constante.

\begin{enumerate}[label=\alph*)]
    \item Explique por qué existe multicolinealidad perfecta.
    \item ¿Qué problema genera para calcular \((X^\top X)^{-1}\)?
    \item ¿Cómo corregiría el problema?
\end{enumerate}

\item \textbf{Discusión aplicada.}

Se quiere estimar el efecto de educación sobre salario. Compare los modelos:
\[
salario_i=\beta_0+\beta_1educ_i+u_i
\]
y
\[
salario_i=\beta_0+\beta_1educ_i+\beta_2exper_i+\beta_3habilidad_i+u_i.
\]

\begin{enumerate}[label=\alph*)]
    \item ¿Por qué el segundo modelo puede entregar una interpretación más cercana a ceteris paribus?
    \item ¿Qué costo estadístico podría tener agregar más variables?
    \item ¿Qué variables sería importante incluir si se busca una interpretación causal?
\end{enumerate}

\end{enumerate}

\section*{Preguntas de cierre}

\begin{enumerate}[label=\arabic*.]
    \item ¿Qué significa mantener constantes las demás variables?
    \item ¿Por qué \(X^\top X\) debe ser invertible?
    \item ¿Qué diferencia hay entre \(R^2\) y \(R^2\) ajustado?
    \item ¿Por qué agregar variables siempre aumenta o mantiene \(R^2\)?
\end{enumerate}

\end{document}
